\newpage

\section{Data Warehouse Architecture}

\subsection{Staging Layer}

\textbf{Input:} fraudTrain.csv (1,296,675 rows) + fraudTest.csv (555,719 rows)

\textbf{Operations:}
\begin{itemize}
    \item Load both CSV files into a single table
    \item Add \texttt{data\_split} column to identify train/test records
    \item Remove duplicate transactions (same transaction\_id)
\end{itemize}

\textbf{Output:} \textbf{fraud\_data} staging table (1,852,394 rows)
\begin{itemize}
    \item Raw transaction data with all original columns intact
    \item Columns: transaction\_id, customer\_id, merchant\_id, amount, transaction\_time, is\_fraud, state, zip, lat, long, city\_pop, merchant\_name, category, data\_split
\end{itemize}

\vspace{0.5cm}

\subsection{Warehouse Layer}

\textbf{Input:} fraud\_data (staging table)

\textbf{Operations:}
\begin{itemize}
    \item Extract unique customers → \textbf{dim\_customer} (999 unique records)
    \item Extract unique merchants → \textbf{dim\_merchant} (693 unique records)
    \item Extract unique locations → \textbf{dim\_location} (985 unique records)
    \item Create hourly time dimension → \textbf{dim\_time} (17,520 hourly buckets)
    \item Denormalize and aggregate transactions → \textbf{fact\_transactions}
\end{itemize}

\textbf{Output Tables:}

\begin{table}[H]
\centering
\begin{tabular}{|l|c|l|}
\hline
\textbf{Table Name} & \textbf{Rows} & \textbf{Purpose} \\
\hline
dim\_customer & 999 & Customer attributes (name, age, gender) \\
dim\_merchant & 693 & Merchant attributes (name, category) \\
dim\_location & 985 & Geographic attributes (state, zip, lat, long, population) \\
dim\_time & 17,520 & Time decomposition (hour, day, month, quarter) \\
fact\_transactions & 1,852,394 & Fact table with transaction details + foreign keys \\
\hline
\end{tabular}
\caption{Warehouse Dimension and Fact Tables}
\label{tab:warehouse_tables}
\end{table}

\textbf{Referential Integrity:} 100\% - All foreign keys in fact\_transactions successfully resolve to dimension tables (0\% NULL values)

\vspace{0.5cm}

\subsection{Data Mart Layer}

\textbf{Input:} Warehouse tables (Dimensions + Fact)

\textbf{Operations:} Create specialized analytical views for different business questions

\textbf{Output Views:}

\begin{enumerate}
    \item \textbf{vw\_fraud\_summary}: Overall fraud statistics
    \begin{itemize}
        \item Total transactions, fraud count, fraud rate
        \item Average transaction amount by fraud status
    \end{itemize}
    
    \item \textbf{vw\_fraud\_by\_merchant}: Merchant-level fraud analysis
    \begin{itemize}
        \item Transactions per merchant, fraud count, fraud rate
        \item Ranking merchants by fraud volume
    \end{itemize}
    
    \item \textbf{vw\_fraud\_by\_location}: Geographic fraud patterns
    \begin{itemize}
        \item Transactions per location, fraud count, fraud rate
        \item Identifying high-risk regions (state, zip)
    \end{itemize}
    
    \item \textbf{vw\_fraud\_by\_customer}: Customer-level fraud analysis
    \begin{itemize}
        \item Customer transaction history, fraud incidents
        \item Identifying repeat fraud victims
    \end{itemize}
    
    \item \textbf{vw\_fraud\_by\_time}: Temporal fraud patterns
    \begin{itemize}
        \item Hourly/daily/monthly fraud trends
        \item Identifying peak fraud times
    \end{itemize}
\end{enumerate}

\vspace{0.5cm}

\subsection{Data Flow Summary}

\begin{table}[H]
\centering
\begin{tabular}{|l|l|l|l|}
\hline
\textbf{Layer} & \textbf{Input} & \textbf{Processing} & \textbf{Output} \\
\hline
\textbf{Staging} & 2 CSV files & Load + Deduplicate & 1 table: fraud\_data \\
& 1.85M rows & & (1.85M rows) \\
\hline
\textbf{Warehouse} & fraud\_data & Dimension extraction & 4 dimensions + 1 fact \\
& & + Aggregation & (1.85M fact records) \\
\hline
\textbf{Data Mart} & Warehouse & Aggregation + Joins & 5 analytical views \\
& & for business queries & (ready for BI tools) \\
\hline
\end{tabular}
\caption{Data Warehouse Architecture Flow}
\label{tab:dataflow}
\end{table}

\section{Data Preprocessing}

\subsection{Overview}

Our preprocessing workflow follows these key steps:
\begin{enumerate}
    \item \textbf{Load data and split}: Load from warehouse and split training set chronologically to simulate real-world conditions
    \item \textbf{Data cleaning}: Validate data quality, confirm no missing values, duplicates, or corruption issues
    \item \textbf{Feature engineering}: Create features from temporal (hour, day, month), geographic (distance), and demographic dimensions (age), and compute behavioral statistics from merchant/cardholder history
    \item \textbf{Feature selection}: Use multiple ranking methods to identify top features with highest discriminative power
    \item \textbf{Data finalization}: Balance training classes with SMOTE and prepare datasets for modeling
\end{enumerate}

\subsection{Data Loading and Splitting}

The first step in preprocessing is loading data from the warehouse fact table and dimension tables. Data splitting strategy is critical in fraud detection because fraud patterns evolve over time.

For this reason, we use \textbf{chronological split} instead of traditional random split. This approach preserves temporal order and ensures the training set always precedes validation/test sets, accurately simulating production conditions where we train on historical data and deploy to predict future transactions.

The raw dataset is split using \textbf{chronological split} to minimize information leakage:
\begin{itemize}
    \item \textbf{Train split}: 80\% of data (1,037,340 rows) allocated for model training
    \item \textbf{Validation split}: 20\% of data (259,335 rows) for hyperparameter tuning, all split chronologically by timestamp
    \item \textbf{Test holdout}: Completely unseen data (separate from train/validation) used only for final evaluation
\end{itemize}

Reasons for using chronological split:
\begin{itemize}
    \item \textbf{Fraud patterns evolve}: Fraud is not a static phenomenon
    \item \textbf{Realistic production simulation}: Train models on transactions from the past and deploy to predict today's transactions
    \item \textbf{Avoid temporal leakage}: Random split might accidentally let validation sets see merchant patterns from future dates
\end{itemize}

\subsection{Data Cleaning and Quality Checks}

Before creating any features, we must ensure the warehouse data is clean and reliable.

\subsubsection{Target Variable Check}

The target column \textit{is\_fraud} is a binary variable (0 = legitimate transaction, 1 = fraud). This dataset is highly imbalanced with fraud comprising only ~0.13\% overall---a major challenge for learning models.

\subsubsection{Duplicate Rows}

Duplicate rows (transactions identical across all columns) are checked and removed. This confirms data integrity.

\subsubsection{Missing Values}

Missing values (NaN, NULL) are checked. Our warehouse data has high quality with minimal missing values.

\subsection{Exploratory Data Analysis (EDA)}

After initial data cleaning, we perform comprehensive exploratory data analysis to understand patterns and distributions that guide feature engineering.

\subsubsection{Key EDA Findings}

From exploratory analysis, several critical fraud patterns emerge:

\textbf{Pattern 1: Temporal Anomalies}
\begin{itemize}
    \item \textbf{Late-night activity}: Fraud spikes between 10 PM - 4 AM
    \item \textbf{Weekend elevation}: Fraud rates are 30-50\% higher on weekends
    \item \textbf{Seasonal patterns}: Holiday periods show elevated fraud
\end{itemize}

\textbf{Pattern 2: Geographic Anomalies}
\begin{itemize}
    \item \textbf{State variation}: Some states show higher fraud rates due to population and transaction volume
    \item \textbf{Distance signal}: Transactions occurring very far from cardholder's typical location are suspicious
    \item \textbf{Impossible travel}: Multiple transactions far apart in short time windows are extreme fraud red flags
\end{itemize}

\textbf{Pattern 3: Amount-Based Signals}
\begin{itemize}
    \item \textbf{Higher amounts}: Fraudulent transactions average \$149 vs \$62 for legitimate (2.4x difference)
    \item \textbf{Outliers}: Very large transactions (>\$1000) warrant additional scrutiny
\end{itemize}

\subsection{Feature Engineering}

Feature engineering creates new features from raw data to capture important patterns. For fraud detection, we create features based on three main sources:
\begin{enumerate}
    \item \textbf{Temporal information}: Transaction timing patterns
    \item \textbf{Geographic information}: Distance from cardholder location to merchant location
    \item \textbf{Behavioral information}: Individual cardholder and merchant patterns
\end{enumerate}

\subsubsection{Temporal Features}

From transaction timestamp, we extract:
\begin{itemize}
    \item \textbf{hour}: Hour of day (0-23)
    \item \textbf{dayofweek}: Day of week (0=Monday, 6=Sunday)
    \item \textbf{month}: Month of year
    \item \textbf{is\_weekend}: Binary variable (1 if weekend)
\end{itemize}

\subsubsection{Geographic Features}

\textbf{distance\_km}: We calculate the Haversine distance from cardholder location to merchant location:
\[\text{distance} = 2R \arcsin\left(\sqrt{\sin^2\left(\frac{\Delta\text{lat}}{2}\right) + \cos(\text{lat}_1) \cos(\text{lat}_2) \sin^2\left(\frac{\Delta\text{lon}}{2}\right)}\right)\]
where $R = 6371$ km is Earth's radius.

\subsubsection{Amount Features}

\textbf{amt\_log}: Log transformation of transaction amount ($\log(1 + \text{amt})$) to reduce outlier impact and normalize distribution.

\subsubsection{Advanced Behavioral Features}

The most important step in fraud detection is capturing behavioral patterns:

\begin{table}[H]
\centering
\caption{Advanced Behavioral Features}
\footnotesize
\begin{tabular}{|l|p{6cm}|}
\hline
\textbf{Feature} & \textbf{Description} \\
\hline
trans\_count\_hourly & Number of cardholder transactions in current hour \\
\hline
trans\_count\_daily & Number of cardholder transactions in current day \\
\hline
merchant\_fraud\_rate & Historical fraud rate of merchant \\
\hline
merchant\_trans\_count & Total transaction count via merchant \\
\hline
merchant\_avg\_amount & Average amount via merchant \\
\hline
amt\_x\_distance & Interaction: amount $\times$ distance \\
\hline
time\_since\_last\_trans & Time elapsed (minutes) since previous transaction \\
\hline
\end{tabular}
\end{table}

\subsection{Feature Scaling and Standardization}

Raw features have dramatically different scales. We apply \textbf{StandardScaler} to transform all features to mean=0 and standard deviation=1:

\[z = \frac{x - \mu}{\sigma}\]

where $\mu$ is the feature mean and $\sigma$ is the standard deviation, both computed from the training set.

\textbf{Critical}: The scaler is fit ONLY on the training set, then applied consistently to all partitions to prevent data leakage.

\subsection{Feature Selection}

At this point, we have ~50-60 features from feature engineering. To select the most valuable features, we use 3 independent ranking methods:

\subsubsection{Correlation-Based Ranking}

Calculate Pearson correlation coefficient between each feature and target variable.

\subsubsection{Mutual Information (MI)}

This method examines dependencies between features and targets, including non-linear relationships:
\[I(X; y) = H(y) - H(y|X)\]

\subsubsection{Tree-Based Feature Importance}

Train a Random Forest and extract feature importance scores from the Gini index.

\subsubsection{Consolidated Score}

Combine all three methods by averaging normalized scores to select top-20 features:
\[Score = \frac{Corr + MI + TreeImp}{3}\]

\subsection{SMOTE Resampling - Balancing Training Classes}

The biggest problem in fraud detection is extreme class imbalance: fraud comprises only 0.166\% of training data. SMOTE (Synthetic Minority Over-sampling Technique) generates synthetic fraud examples:
\begin{itemize}
    \item Find k-nearest neighbors of each fraud example
    \item Create synthetic points between fraud examples and neighbors
    \item Repeat until balanced (fraud = legitimate = 50\%)
\end{itemize}

Result: balanced training set helps models learn patterns of both classes. However, validation/test are kept original (not balanced) to reflect reality.

\begin{table}[H]
\centering
\caption{Final 20 Features for Modeling}
\small
\begin{tabular}{|c|l|}
\hline
\textbf{No} & \textbf{Feature} \\
\hline
1 & distance\_km \\
\hline
2 & amt\_log \\
\hline
3 & merchant\_fraud\_rate \\
\hline
4 & hour \\
\hline
5 & trans\_count\_hourly \\
\hline
6 & trans\_count\_daily \\
\hline
7 & amt \\
\hline
8 & merchant\_avg\_amount \\
\hline
9 & age \\
\hline
10 & is\_weekend \\
\hline
11 & dayofweek \\
\hline
12 & category \\
\hline
13 & month \\
\hline
14 & amt\_x\_distance \\
\hline
15 & merchant\_trans\_count \\
\hline
16 & time\_since\_last\_trans \\
\hline
17 & age\_group \\
\hline
18 & state \\
\hline
19 & gender \\
\hline
20 & fraud\_dist\_percentile \\
\hline
\end{tabular}
\end{table}

\subsection{Data Ready for Modeling}

After all preprocessing and feature selection steps, data is ready for modeling with 3 separate datasets:

\begin{itemize}
    \item \textbf{Training set} (1,037,340 rows): Used to train the model, SMOTE-balanced
    \item \textbf{Validation set} (259,335 rows): Used to tune hyperparameters, original imbalanced distribution
    \item \textbf{Test holdout} (555,719 rows): Final evaluation set, completely unseen data
\end{itemize}

\begin{table}[H]
\centering
\caption{Summary of Final Data for Modeling}
\begin{tabular}{|c|c|c|c|c|}
\hline
\textbf{Dataset} & \textbf{Rows} & \textbf{Features} & \textbf{Fraud Rate} & \textbf{Purpose} \\
\hline
Training & 1,037,340 & 20 & 0.166\% & Train the model \\
\hline
Validation & 259,335 & 20 & 0.089\% & Tune hyperparameters \\
\hline
Test holdout & 555,719 & 20 & 0.080\% & Final evaluation \\
\hline
\end{tabular}
\end{table}

